% !TeX root = ../../Book1.tex
% 纲要标题：### 11.4 浸入、淹没与嵌入
% 纲要位置：几何学计划纲要.md，第 558 行。
% 正文状态：初稿；尚未正式审核。

% 纲要内容开始
% 
% * 浸入和淹没的标准形
% * 正则水平集定理
% * 嵌入、初始子流形结构及限制映射
% * 单射浸入未必嵌入
% * 带边界版本中切空间与边界条件
% 
% ---
% 
% 纲要内容结束

\section{浸入、淹没与嵌入}
\label{g1:sec:1104}

第八章用秩为二的参数映射描述曲面，现在可以把同一个问题放到任意维数的流形上。一张参数片是否丢失切向方向，由微分是否单射判断；一个方程组能否在附近独立控制所有输出，由微分是否满射判断。但即使每张参数片都很规整，远处的参数仍可能回到眼前，这就需要另加拓扑条件。常秩定理先给出局部图像，嵌入的定义再保证这些局部图像与整个像的拓扑相符。本节前四个小节先讨论无边界流形，最后专门说明边界处哪些条件需要补充。

\subsection{浸入和淹没的标准形}
\label{g1:subsec:1104:01}

\begin{definition}[浸入与淹没]\label{g1:def:immersion-submersion}
	设 \(F:M^m\to N^n\) 是光滑映射。若 \(\mathrm{d}F_p:\mathrm{T}_pM\to \mathrm{T}_{F(p)}N\) 单射，称 \(F\) 在 \(p\) 处为浸入；若它满射，称 \(F\) 在 \(p\) 处为\term[yanmo]{淹没}[submersion]。处处满足相应条件时，分别称 \(F\) 为浸入或淹没。
\end{definition}

这里的浸入推广了\cref{g1:def:surface-immersion}中的曲面参数：其定义域维数为二，所以满列秩正是两个参数方向不被压在一起。一般浸入必须有 \(m\le n\)，一般淹没必须有 \(m\ge n\)。淹没也不要求映射在集合意义下满射，例如 \(t\mapsto e^t\) 是 \(\mathbb R\to\mathbb R\) 的淹没，而像只是正半轴。

\begin{theorem}[浸入与淹没的局部标准形]\label{g1:thm:immersion-submersion-normal-form}
	设 \(M^m,N^n\) 是无边界光滑流形，\(F:M\to N\) 光滑，\(p\in M\)。
	\begin{enumerate}
		\item 若 \(\mathrm{d}F_p\) 单射，则存在分别以 \(p,F(p)\) 为原点的源、目标坐标，使 \(F\) 在源的一个邻域中写成
		\[
			(u^1,\ldots,u^m)\longmapsto(u^1,\ldots,u^m,0,\ldots,0).
		\]
		\item 若 \(\mathrm{d}F_p\) 满射，则可选这样的坐标，使 \(F\) 写成
		\[
			(u^1,\ldots,u^m)\longmapsto(u^1,\ldots,u^n).
		\]
	\end{enumerate}
\end{theorem}
\begin{proof}
	单射情形，坐标 Jacobian 在 \(p\) 处有一个非零的 \(m\) 阶子式，连续性使它在某邻域内仍非零。秩又不能超过源维数 \(m\)，故该邻域的秩恒为 \(m\)。满射情形同理：一个非零的 \(n\) 阶子式保证邻域内的秩恒为 \(n\)。现在分别对秩 \(m\) 和秩 \(n\) 使用\cref{g1:cor:manifold-constant-rank}，就得到两种标准形。这里先证明邻域内秩恒定，才能调用常秩结论；仅知道一点的非极大秩，通常不能这样做。
\end{proof}

从标准形可以读出两项常用结论。浸入限制到足够小的源邻域后，是到其像的同胚：在坐标中只需投影到前 \(m\) 个坐标就得到连续逆映射。这说的是一个指定的参数分支，并没有断言其他参数分支不会进入同一个目标邻域。

\ProseParagraphBreak

无边界流形之间的淹没是开映射。事实上，把任一源开集中的点放进包含在该开集内的坐标小盒，投影标准形把这个小盒送到目标中的开盒，所有这些像的并就是原开集的像。此外，若 \(q=F(p)\)，在投影标准形中取被忘掉的坐标恒为零，就得到一个通过 \(p\) 的光滑局部截面 \(s\)，满足 \(F\circ s=\operatorname{id}\)。因此淹没允许在目标点附近光滑选择一个原像，但未必能把这些局部选择拼成全局截面。

\subsection{正则水平集定理}
\label{g1:subsec:1104:02}

方程 \(F(p)=y\) 应当减少多少自由度，取决于它实际提供了多少个独立约束。若目标维数为 \(n\)，微分满射正是所有 \(n\) 个约束在一阶上独立。淹没标准形将这个判断直接变成坐标中固定前 \(n\) 个数。

\begin{definition}[正则值]\label{g1:def:regular-value}
	设 \(F:M\to N\) 光滑。若 \(\mathrm{d}F_p\) 满射，称 \(p\) 为 \(F\) 的\term[zhengzedian]{正则点}[regular point]；否则称为\term[linjiedian]{临界点}[critical point]。若 \(y\in N\) 的每个原像都是正则点，称 \(y\) 为\term[zhengzezhi]{正则值}[regular value]；否则称为\term[linjiezhi]{临界值}[critical value]。没有原像的点也按此定义算作正则值。
\end{definition}

\begin{theorem}[正则水平集]\label{g1:thm:regular-level-set}
	设 \(F:M^m\to N^n\) 是无边界光滑流形之间的光滑映射，\(y\in N\) 是正则值。若 \(S=F^{-1}(y)\) 非空，则 \(S\) 在子空间拓扑下有自然的 \((m-n)\) 维光滑流形结构，包含映射 \(i:S\hookrightarrow M\) 是浸入，且
	\begin{equation}\label{g1:eq:regular-level-tangent}
		\mathrm{d}i_p(\mathrm{T}_pS)=\ker\mathrm{d}F_p\qquad(p\in S).
	\end{equation}
	通过单射 \(\mathrm{d}i_p\) 把 \(\mathrm{T}_pS\) 视为 \(\mathrm{T}_pM\) 的子空间时，可简写为 \(\mathrm{T}_pS=\ker\mathrm{d}F_p\)。
\end{theorem}
\begin{proof}
	对每个 \(p\in S\) 使用淹没标准形，将 \(y\) 的目标坐标取为零。在所得源坐标 \((u^1,\ldots,u^m)\) 中，水平集恰为
	\[
		u^1=\cdots=u^n=0.
	\]
	因此后 \(m-n\) 个坐标给出 \(S\) 的局部坐标，而且这些坐标与 \(S\) 的子空间拓扑一致。两张这样的坐标之间，转换映射是在环境坐标转换中先代入前 \(n\) 个坐标为零，再取所需坐标分量，故光滑；逆转换同理。这构成光滑图册。Hausdorff 性和第二可数性从环境的子空间继承，所以确实得到流形。

	在这些坐标中，包含映射的微分把后 \(m-n\) 个坐标方向单射送入 \(\mathrm{T}_pM\)。另一方面，\(\mathrm{d}F_p\) 是前 \(n\) 个方向的投影，其核恰由后 \(m-n\) 个方向张成。于是得到\eqref{g1:eq:regular-level-tangent}。坐标改变通过\cref{g1:prop:differential-chain-rule}的微分链式法则相容，所以这个子空间描述不依赖所选坐标。
\end{proof}

下一小节将通过包含映射的性质进一步说明这个自然结构为何唯一。这里并非随意给同一集合指定一种流形拓扑。只需在水平集各点满足正则值的条件，不要求 \(F\) 在整个 \(M\) 上都是淹没。

\begin{example}\label{g1:exm:regular-sphere-and-singular-level}
	设 \(F:\mathbb R^{k+1}\to\mathbb R\)，\(F(x)=\|x\|^2\)，其中 \(k\ge1\)。因为
	\[
		\mathrm{d}F_x(v)=2\langle x,v\rangle,
	\]
	每个 \(c>0\) 都是正则值，球面 \(F^{-1}(c)\) 的维数为 \(k\)，并且 \(\mathrm{T}_xF^{-1}(c)=x^\perp\)。负数也为正则值，但水平集为空。零为临界值，其水平集只有一个点；该集合仍可以是零维流形，却不具有正则水平集定理所给的预期维数 \(k\)。

	另取 \(G(x,y)=xy\)。它的零水平集是两条坐标轴的并，在原点形成四条分支，不是原点附近的一维流形。两个例子说明：临界水平集有时仍规整，有时出现交叉，仅由“不是正则值”不能断言所有水平集具有同一种奇异性。
\end{example}

\subsection{嵌入、初始子流形结构及限制映射}
\label{g1:subsec:1104:03}

第四章的拓扑嵌入要求映射与其像之间互为连续逆映射，像必须使用环境的子空间拓扑。加入这个条件以后，浸入的局部参数片才会与整个像的局部邻域一致。

\begin{definition}[光滑嵌入]\label{g1:def:smooth-embedding}
	无边界光滑流形之间的光滑映射 \(f:M\to N\)，若既是浸入，又是 \(M\) 到 \(f(M)\) 的同胚，其中 \(f(M)\) 取 \(N\) 的子空间拓扑，则称为\term[guanghuaqianru]{光滑嵌入}[smooth embedding]，简称嵌入。
\end{definition}

\begin{definition}[嵌入子流形]\label{g1:def:embedded-submanifold}
	设 \(N^n\) 为无边界光滑流形。子集 \(S\subset N\) 若对每个 \(p\in S\) 都有环境坐标 \((V,y)\)，使
	\[
		y(S\cap V)=y(V)\cap\bigl(\mathbb R^k\times\{0\}^{n-k}\bigr),
	\]
	则称为 \(N\) 的 \(k\) 维\term[qianruziliuxing]{嵌入子流形}[embedded submanifold]。给 \(S\) 子空间拓扑，并用这些切片的前 \(k\) 个坐标作为图册。
\end{definition}

坐标切片之间的转换正如正则水平集证明中那样光滑，所以这个定义确实给出 \(S\) 的光滑结构。下述等价关系说明，这个集合式定义与前面的映射式定义对应。

\begin{proposition}[嵌入与坐标切片]\label{g1:prop:embedding-slice-characterization}
	设 \(M^k,N^n\) 是无边界光滑流形。若 \(f:M\to N\) 为光滑嵌入，则其像 \(S=f(M)\) 为嵌入子流形，且 \(f:M\to S\) 是微分同胚。反之，每个嵌入子流形的包含映射都是光滑嵌入。
\end{proposition}
\begin{proof}
	取 \(p\in M\)，浸入标准形给出源邻域 \(U\) 及 \(f(p)\) 附近的坐标，使 \(f|_U\) 的像是一块坐标切片。必要时缩小坐标盒，可以使这块切片完整覆盖目标小盒中的前 \(k\) 个坐标。关键是排除 \(M\setminus U\) 的点也映入眼前的小盒：因为 \(f\) 是到子空间 \(S\) 的同胚，\(f(U)\) 在 \(S\) 中开，故存在环境开集 \(W\subset N\)，满足 \(S\cap W=f(U)\)。把目标坐标盒进一步缩到 \(W\) 内，盒内的整个 \(S\) 就恰为该切片，而不只是其中一个参数分支。这样的盒覆盖 \(S\)，给出嵌入子流形结构。在对应坐标中，\(f:M\to S\) 是恒等坐标映射，故连同逆映射都光滑。

	反过来，包含映射在适配坐标中为 \(u\mapsto(u,0)\)，所以是浸入。\(S\) 已取环境的子空间拓扑，包含映射自然是到其像的同胚，故为光滑嵌入。
\end{proof}

\begin{proposition}[嵌入子流形的初始性质]\label{g1:prop:embedded-initial-property}
	设 \(S\subset N\) 为无边界光滑流形 \(N\) 的嵌入子流形，\(i:S\hookrightarrow N\) 为包含映射。对任意光滑流形 \(P\)（允许有边界）及任意集合映射 \(h:P\to S\)，有
	\[
		h:P\to S\text{ 光滑}\quad\Longleftrightarrow\quad i\circ h:P\to N\text{ 光滑}.
	\]
	特别地，光滑映射 \(H:P\to N\) 若像包含在 \(S\) 中，就唯一地给出光滑映射 \(P\to S\)。在固定子集 \(S\subset N\) 上，使其成为嵌入子流形的光滑结构唯一。
\end{proposition}
\begin{proof}
	正向由光滑映射的复合得到。反向先处理一个容易遗漏的拓扑问题：若 \(i\circ h\) 光滑，则它连续；由于 \(S\) 取子空间拓扑，\(h\) 也连续。于是对 \(h(p)\) 附近的适配图 \(V\)，集合 \(h^{-1}(S\cap V)\) 在 \(P\) 中开。在这片上，\(h\) 的坐标表示就是 \(i\circ h\) 的坐标表示的前 \(\dim S\) 个分量，所以光滑。即使 \(p\) 在 \(\partial P\)，这些分量仍由环境中的光滑延拓限制得到，论证不变。局部结论合起来给出 \(h\) 光滑。

	像包含在 \(S\) 的映射只有一种集合意义下的限制值域方式，上述等价立即证明其光滑性。若同一个子集带有两种嵌入子流形结构，两种包含映射都光滑。对两种结构之间的恒等集合映射分别使用已证等价，得到恒等映射在两个方向都光滑，故两种极大光滑图册相同。
\end{proof}

这里的“初始”说的是通过包含映射检测光滑性的性质，不是另给 \(S\) 加一种比子空间拓扑更细的拓扑。例如光滑函数 \(f:N\to\mathbb R\) 限制到 \(S\) 后光滑；若 \(f(S)\subset T\)，其中 \(T\) 也是嵌入子流形，则 \(f|_S:S\to T\) 光滑。正则水平集由环境坐标切片构造，因而属于这一情形。

\ProseParagraphBreak

包含映射的微分通常被省略在记号里：对嵌入子流形写 \(\mathrm{T}_pS\subset \mathrm{T}_pN\)，实际是用单射 \(\mathrm{d}i_p\) 作识别。这不是把两个定义不同的导子集合逐项当作相同集合，而是利用包含映射，在前者与后者的一个子空间之间建立自然的线性同构。

\subsection{单射浸入未必嵌入}
\label{g1:subsec:1104:04}

单射解决了“一个像点有几个原像”，却没有解决“附近像点的原像是否也在附近”。下面的平面曲线已经能把两者区别开，不需要依靠高维或稠密轨道的例子。

\begin{example}\label{g1:exm:injective-immersion-not-embedding}
	令
	\[
		\gamma:(-\pi,\pi)\longrightarrow\mathbb R^2,
		\qquad\gamma(t)=(\sin t,\sin2t).
	\]
	它是单射浸入，但不是嵌入。先检查浸入：\(\gamma'(t)=(\cos t,2\cos2t)\)，若第一分量为零，则 \(t=\pm\pi/2\)，第二分量为 \(-2\ne0\)，故导数从不为零。

	再检查单射。\(\sin t\) 的正负区分正、负参数，且在本区间内只在 \(t=0\) 为零。若 \(s,t\in(0,\pi)\) 且 \(\sin s=\sin t\)，则 \(s=t\) 或 \(s=\pi-t\)；后一情形的 \(\sin2s=-\sin2t\)，两者相等只能有 \(t=\pi/2\)，此时仍为 \(s=t\)。负参数的情形用 \(\gamma(-t)=-\gamma(t)\) 化为正参数，因而 \(\gamma\) 单射。

	然而取 \(t_j\to\pi\)，便有 \(\gamma(t_j)\to(0,0)=\gamma(0)\)，而 \(t_j\) 不趋于零。因此从像的子空间拓扑返回参数的逆映射不连续。\cref{g1:fig:injective-immersion-topology}中，蓝色短弧来自零附近的参数，另一方向逼近交点的曲线来自接近参数区间端点的部分。浸入标准形只能控制前一支。

	这个像甚至是闭集：把同一公式延拓到紧区间 \([-\pi,\pi]\)，两个端点都映到已有的 \(\gamma(0)\)，所以延拓前后的像相同，是紧集。可见“单射浸入且像闭”仍不能替代嵌入条件。
\end{example}

% !TeX root = ../Book1.tex
\begin{figure}[htbp]
	\centering
	\begin{tikzpicture}[font=\small]
		\begin{scope}[xshift=-2.8cm,x=2cm,y=1.3cm]
			\draw[->,GeometryGray] (-1.25,0)--(1.3,0) node[right] {\(x\)};
			\draw[->,GeometryGray] (0,-1.2)--(0,1.25) node[above] {\(y\)};
			\draw[GeometryGray,thick,domain=-180:180,samples=181,variable=\t] plot ({sin(\t)},{sin(2*\t)});
			\draw[GeometryBlue,very thick,domain=-27:27,samples=40,variable=\t] plot ({sin(\t)},{sin(2*\t)});
			\draw[orange!75!black,very thick,domain=145:179.8,samples=35,variable=\t] plot ({sin(\t)},{sin(2*\t)});
			\foreach \t in {150,165,174,178}{\fill[orange!75!black] ({sin(\t)},{sin(2*\t)}) circle (1.25pt);}
			\fill[GeometryBlue] (0,0) circle (1.6pt);
			\node[above left] at (0,0) {\(\gamma(0)\)};
			\node[GeometryBlue,above right] at (0.5,0.78) {\(t\approx0\)};
			\node[orange!75!black,right] at (0.63,-0.72) {\(t\uparrow\pi\)};
		\end{scope}
		\begin{scope}[xshift=3.7cm,yshift=0.2cm]
			\draw[GeometryGray] (-1.7,0)--(1.7,0);
			\draw[fill=white] (-1.7,0) circle (2pt);
			\draw[fill=white] (1.7,0) circle (2pt);
			\fill[GeometryBlue] (0,0) circle (2pt);
			\node[below] at (-1.7,0) {\(-\pi\)};
			\node[below] at (0,0) {\(0\)};
			\node[below] at (1.7,0) {\(\pi\)};
			\foreach \x in {1,1.35,1.52,1.61}{\fill[orange!75!black] (\x,0) circle (1.5pt);}
			\draw[->,orange!75!black] (0.95,0.35)--(1.65,0.35);
			\node[above] at (0,0.75) {参数区间 \((-\pi,\pi)\)};
			\node[align=center] at (0,-1) {\(\gamma(t_j)\to\gamma(0)\)\\但 \(t_j\not\to0\)};
		\end{scope}
	\end{tikzpicture}
	\caption[单射浸入的像拓扑]{曲线 \(\gamma(t)=(\sin t,\sin2t)\)，定义域为 \((-\pi,\pi)\)。蓝色弧是零附近的参数片；橙色点来自趋向 \(\pi\) 的参数，也在像中趋向原点。逆参数映射因而不连续，尽管每个参数点的导数非零且参数映射单射。}
	\label{g1:fig:injective-immersion-topology}
\end{figure}


还可以在这个例子中具体看出，嵌入子流形的初始性质不能自动移植给任意单射浸入像。给集合 \(S=\gamma((-\pi,\pi))\) 从参数区间搬来流形结构，使 \(\gamma\) 成为微分同胚；此时包含 \(S\to\mathbb R^2\) 是浸入。对充分小的 \(\varepsilon>0\)，考虑值在 \(S\) 中的曲线
\[
	h(s)=(-\sin s,\sin2s),\qquad |s|<\varepsilon.
\]
作为到 \(\mathbb R^2\) 的映射，它光滑；但它在 \(S\) 的参数坐标为
\[
	\gamma^{-1}(h(s))=
	\begin{cases}
		\pi+s,&s<0,\\
		0,&s=0,\\
		-\pi+s,&s>0,
	\end{cases}
\]
在零处不连续，故 \(h:(-\varepsilon,\varepsilon)\to S\) 不是光滑映射。问题恰是这份流形结构的拓扑不同于 \(S\) 的子空间拓扑。

\begin{definition}[固有映射]\label{g1:def:proper-map-manifolds}
	光滑流形之间的连续映射 \(f:M\to N\) 若对每个紧集 \(K\subset N\)，原像 \(f^{-1}(K)\) 都紧，则称为\term[guyouyingshe]{固有映射}[proper map]。
\end{definition}

固有性控制的是“落在一个固定紧目标区域内的所有参数”。上一例中整个像紧，而其原像是非紧区间 \((-\pi,\pi)\)，所以 \(\gamma\) 不固有。这一条件能排除远处参数悄悄回到紧区域的现象。

\begin{proposition}[固有单射浸入]\label{g1:prop:proper-injective-immersion}
	设 \(M,N\) 为无边界光滑流形，\(f:M\to N\) 为单射浸入。若 \(f\) 固有，则 \(f\) 是嵌入，且像闭于 \(N\)。特别地，若 \(M\) 紧，则每个单射浸入 \(M\to N\) 都是闭像嵌入。反之，闭像嵌入是固有映射。
\end{proposition}
\begin{proof}
	先证明一个连续固有映射 \(f:M\to N\) 是闭映射。取闭集 \(A\subset M\) 及 \(y\in\overline{f(A)}\)。在 \(N\) 的坐标图内取一个紧邻域 \(K\ni y\)，使 \(y\in\operatorname{int}K\)。固有性使 \(f^{-1}(K)\) 紧，故其闭子集 \(A\cap f^{-1}(K)\) 紧。它的连续像
	\[
		f(A\cap f^{-1}(K))=f(A)\cap K
	\]
	也紧，并由\cref{g1:prop:compact-hausdorff-closed}闭于 \(N\)。回忆这里闭性的理由：对紧集外一点，用 Hausdorff 性把它与紧集的每一点用开邻域分开，再以有限子覆盖替代紧集一侧的全部邻域，取外点一侧有限个邻域的交，便得到与紧集不交的开邻域。因为 \(y\) 位于 \(K\) 的内部且是 \(f(A)\) 的闭包点，它也是 \(f(A)\cap K\) 的闭包点，因而属于该闭集。于是 \(f(A)\) 闭。

	现在 \(f\) 连续且单射，并将闭集送到闭集，所以 \(f:M\to f(M)\) 的逆映射连续。它又是浸入，故为嵌入；取 \(A=M\) 还得到像闭。若 \(M\) 紧，每个目标紧集在 Hausdorff 空间 \(N\) 中闭，其原像为 \(M\) 的闭子集，因而紧，所以 \(f\) 自动固有。此时也可直接使用\cref{g1:thm:compact-bijection}判断到像的逆连续。

	反之，若 \(f\) 是闭像嵌入，则 \(K\cap f(M)\) 对每个紧 \(K\subset N\) 都是紧集；由同胚 \(M\cong f(M)\)，其原像 \(f^{-1}(K)\) 紧，故 \(f\) 固有。
\end{proof}

并非所有嵌入都固有。例如开区间 \((0,1)\hookrightarrow\mathbb R\) 是嵌入，但紧集 \([0,1]\) 的原像不紧。固有单射浸入给出的是一种保证闭像嵌入的充分判据，并没有把“嵌入”的含义改成“闭子集”。

\subsection{带边界版本中切空间与边界条件}
\label{g1:subsec:1104:05}

对带边界流形，浸入、淹没仍按完整微分的单射、满射定义，嵌入仍要求同时是浸入和到子空间像的同胚。边界坐标采用 \(\mathbb H^m=\{x^m\ge0\}\)。\cref{g1:prop:derivation-coordinate-basis}及\cref{g1:prop:boundary-velocities}说明，即使 \(p\in\partial M\)，完整切空间 \(\mathrm{T}_pM\) 仍是 \(m\) 维向量空间；向内半空间只是其中能够由向内单侧曲线实现的速度集合，不能代替整个切空间。若以 \(\rho=x^m\) 为局部边界定义函数，则
\[
	\mathrm{T}_p(\partial M)=\ker\mathrm{d}\rho_p\subset \mathrm{T}_pM.
\]
边界上各种满秩条件因此仍按完整切空间计算，但由此得到的 Euclid 局部变换未必保留半空间。

\begin{proposition}[边界处的局部逆]\label{g1:prop:boundary-local-inverse}
	设 \(M^m,N^m\) 为带边界光滑流形，\(m\ge1\)，\(F:M\to N\) 光滑。若 \(p\in\partial M\)、\(F(p)\in\partial N\)，在 \(p\) 的某邻域内有 \(F(\partial M)\subset\partial N\)，并且 \(\mathrm{d}F_p\) 可逆，则 \(F\) 限制到 \(p\) 的某个相对开邻域后，是到 \(F(p)\) 的一个相对开邻域的微分同胚。
\end{proposition}
\begin{proof}
	取两边的半空间坐标，把 \(p,F(p)\) 都置于零。按边界光滑映射的定义，坐标表示 \(G\) 可光滑延拓到 \(\mathbb R^m\) 的一个开邻域。保边界条件给出 \(G^m(x',0)=0\)，所以由一元微积分基本定理，
	\[
		G^m(x',x^m)=x^m a(x',x^m),\qquad
		a(x',x^m)=\int_0^1\partial_mG^m(x',t x^m)\,\mathrm{d}t.
	\]
	\(\mathrm{d}G_0\) 的最后一行在前 \(m-1\) 列为零；可逆性因而使 \(a(0)\ne0\)。又因为 \(G\) 把 \(x^m\ge0\) 的点送到目标半空间，一侧导数给出 \(a(0)\ge0\)，故 \(a(0)>0\)。缩小邻域使 \(a>0\)，则 \(G^m\) 与 \(x^m\) 同号。对 Euclid 延拓使用\cref{g1:thm:inverse-function}，得到开邻域间的光滑逆；同号性质说明该逆也保留半空间。把这对互逆映射限制到半空间，便得到所需边界局部逆。
\end{proof}

保边界条件不能只检查在一个点成立。例如
\[
	G:\mathbb H^2\to\mathbb H^2,\qquad G(x,y)=(x,y+x^2)
\]
在原点的微分是恒等映射，也把原点送到原点，但它把附近的边界点送到抛物线 \(v=u^2\)。其像在 \(v\ge u^2\) 中，不包含原点在目标半平面中的任何相对开邻域，所以不是边界局部微分同胚。Euclid 延拓的逆 \((u,v)\mapsto(u,v-u^2)\) 光滑，却未必取值于源半平面。

\ProseParagraphBreak

对水平集，正确的附加条件要求约束在边界切向方向中仍然独立。这样，约束不会把唯一的向内方向全部吃掉。若带边界子流形 \(S\subset M\) 满足 \(\partial S=S\cap\partial M\)，且在这些点满足
\[
	\mathrm{T}_pS+\mathrm{T}_p(\partial M)=\mathrm{T}_pM,
\]
本书用\term[neatziliuxing]{与边界相容的横截嵌入子流形}[neat submanifold]描述同时满足这两个条件的子流形。这里包含映射要求为嵌入；横截等式意味着 \(\mathrm{T}_pS\) 中还保留至少一个不沿环境边界的方向。

\begin{theorem}[带边界的正则水平集]\label{g1:thm:boundary-regular-level-set}
	设 \(M^m\) 为带边界光滑流形，\(N^n\) 为无边界光滑流形，\(F:M\to N\) 光滑。若 \(y\in N\) 同时为 \(F\) 及 \(F|_{\partial M}\) 的正则值，则非空水平集 \(S=F^{-1}(y)\) 自然成为 \((m-n)\) 维带边界嵌入子流形，且
	\[
		\partial S=S\cap\partial M,\qquad \mathrm{T}_pS=\ker\mathrm{d}F_p.
	\]
	它是与边界相容的横截嵌入子流形；在边界点还满足
	\[
		\mathrm{T}_p(\partial S)=\ker\mathrm{d}F_p\cap \mathrm{T}_p(\partial M).
	\]
\end{theorem}
\begin{proof}
	内部点使用\cref{g1:thm:regular-level-set}的局部证明。取 \(p\in S\cap\partial M\)，用半空间坐标将 \(p\) 置于零，用目标坐标将 \(y\) 置于零，记映射的坐标分量为 \(f^1,\ldots,f^n\)。\(F|_{\partial M}\) 在 \(p\) 处满秩，意味着 Jacobian 的前 \(m-1\) 列中有一个非零的 \(n\) 阶子式。重排前 \(m-1\) 个坐标，可令该子式来自前 \(n\) 列。特别地，若水平集碰到边界，必有 \(m-n\ge1\)。

	对 \(f\) 的 Euclid 光滑延拓构造
	\[
		\Phi(x)=\bigl(f^1(x),\ldots,f^n(x),x^{n+1},\ldots,x^{m-1},x^m\bigr).
	\]
	\(\mathrm{d}\Phi_0\) 可逆，故由\cref{g1:thm:inverse-function}，它在足够小的 Euclid 开邻域间为微分同胚。与任意使用常秩定理不同，这里特意把最后一个坐标原样保留：\(\Phi^m(x)=x^m\)。因此 \(\Phi\) 及其逆都保留半空间和边界。以 \(u=\Phi(x)\) 为新的边界坐标，水平集成为
	\begin{equation}\label{g1:eq:boundary-level-slice}
		u^1=\cdots=u^n=0,\qquad u^m\ge0.
	\end{equation}
	后 \(m-n\) 个坐标给出水平集的半空间图，其边界恰是 \(u^m=0\)，也就是环境边界与水平集的交。图册转换由环境光滑转换的限制得到，子空间拓扑、Hausdorff 性及第二可数性同无边界情形，因此得到带边界嵌入子流形。

	在这些坐标中，\(\mathrm{T}_pS\) 由后 \(m-n\) 个完整坐标方向张成，等于 \(\ker\mathrm{d}F_p\)；其边界切空间再去掉最后一个方向，正是 \(\ker\mathrm{d}F_p\cap \mathrm{T}_p(\partial M)\)。最后一个方向 \(\partial/\partial u^m\) 属于 \(\mathrm{T}_pS\) 却不属于 \(\mathrm{T}_p(\partial M)\)，而后者余维为一，故两者之和为整个 \(\mathrm{T}_pM\)。这就同时验证了边界交等式与横截条件。
\end{proof}

上述半空间切片还说明，与边界相容的横截水平集仍具有初始性质：一个映射若值在 \(S\) 中，其是否光滑可以通过组合包含 \(S\hookrightarrow M\) 来检验。证明仍先用子空间拓扑得到连续性，再取\eqref{g1:eq:boundary-level-slice}中后 \(m-n\) 个坐标；这些坐标有环境光滑延拓，因此同样适用于带边界的源。

\ProseParagraphBreak

只检查完整切空间上的 \(\mathrm{d}F\) 满射还不够。例如 \(F:\mathbb H^m\to\mathbb R\)，\(F(x)=x^m\)，处处满秩，但 \(F^{-1}(0)=\partial M\) 是一张无边界的 \((m-1)\) 维流形；其自身边界为空，而与 \(\partial M\) 的交却是整个水平集，故不是 与边界相容的横截水平集。更明显地，\(F(x,y)=y+x^2\) 在 \(\mathbb H^2\) 上也处处满秩，但零水平集只有原点，连预期的一维局部形状都没有。两例中 \(F|_{\partial M}\) 都在相关点失去满秩性，正是附加条件所要排除的现象。
